\begin{abstract}
Vision--language models (VLMs) are increasingly deployed for video
understanding, yet their ability to reason about \emph{who} did
\emph{what} to \emph{whom} in aggressive interactions remains largely
unexamined.
We introduce a benchmark of 2,687 video clips spanning 20 aggression
categories, accompanied by over 16,500 multiple-choice questions
organized into four tiers of compositional difficulty, from single-attribute
identification to full aggressor--action--victim reasoning.
Each question's distractors are constructed from a 10-level hardness
taxonomy that systematically controls for shortcut cues such as role
reversal, bystander substitution, and cross-video borrowing.
Zero-shot evaluation of 11 VLMs (7--72B parameters) reveals that even
the strongest model achieves only 46\% accuracy on primary questions,
with performance dropping below 34\% on compositional questions
requiring joint aggressor--action--victim understanding.
To address this gap, we propose a progressive fine-tuning pipeline:
supervised fine-tuning teaches MCQ format compliance, chain-of-thought
distillation adds structured reasoning for compositional questions, and
anchored direct preference optimization sharpens discrimination against
taxonomy-hard distractors.
Training on just 20\% of videos with group-aware splits that prevent
parent-video leakage, our pipeline improves accuracy to \textbf{XX\%}
(+XX points), with the largest gains on the compositional categories
that zero-shot models struggle with most.
\end{abstract}


Vision--language models (VLMs) are increasingly deployed for video understanding, yet their ability to reason about \emph{who} did \emph{what} to \emph{whom} in aggressive interactions remains largely unexamined.We introduce a benchmark of 2,687 video clips spanning 20 aggression categories, accompanied by over 16,500 multiple-choice questions organized into four tiers of compositional difficulty, from single-attribute identification to full aggressor--action--victim reasoning. Each question's distractors are constructed from a 10-level hardness taxonomy that systematically controls for shortcut cues such as role reversal, bystander substitution, and cross-video borrowing. Zero-shot evaluation of 11 VLMs reveals that even the strongest model achieves only XX\% accuracy on primary questions, with performance dropping below XX\% on compositional questions requiring joint aggressor--action--victim understanding.
  
To address this gap, we propose a progressive fine-tuning pipeline: supervised fine-tuning teaches MCQ format compliance, chain-of-thought distillation adds structured reasoning for compositional questions, and anchored direct preference optimization sharpens discrimination against taxonomy-hard distractors.
Training on just 20\% of videos with group-aware splits that prevent
parent-video leakage, our pipeline improves accuracy to \textbf{XX\%}
(+XX points), with the largest gains on the compositional categories
that zero-shot models struggle with most.